\begin{paracol}{2}
\section*{QA2b. Theoretical Q\&A — Simplification \& Denoising}

\textbf{Q1. What is mesh simplification? Why is it needed?}\\
\textbf{Mesh simplification} is the process of reducing the number of vertices, edges, and faces in a triangle mesh while preserving its visual and geometric appearance as much as possible. Modern scanners and CAD tools easily produce meshes with millions of triangles, which are too heavy for real-time rendering, transmission, or physics simulation. Simplification creates lightweight models suitable for \textbf{LOD (level of detail)}, mobile/VR devices, and web streaming. It also removes redundant micro-detail that does not contribute to perception. Algorithms range from \textbf{vertex clustering}, \textbf{vertex decimation} (Schroeder), to \textbf{QEM edge collapse} (Garland \& Heckbert). The goal is to minimize geometric error subject to a target triangle budget.

\textbf{Q2. Explain Quadric Error Metric (QEM) edge collapse — write the math.}\\
For each vertex $v$, define the \textbf{quadric} $Q = \sum_{p \in \text{planes}(v)} K_p$, where each plane $p: ax+by+cz+d=0$ contributes $K_p = \mathbf{p}\mathbf{p}^T$ with $\mathbf{p}=(a,b,c,d)^T$. The error of placing $v=(x,y,z,1)^T$ is the sum of squared distances to those planes: $\Delta(v) = v^T Q v$. For an \textbf{edge collapse} $(v_i,v_j) \to \bar v$, the merged quadric is $\bar Q = Q_i + Q_j$ and the error is $\Delta(\bar v) = \bar v^T \bar Q \bar v$. Edges are inserted into a heap keyed by this minimum error and collapsed in order. After each collapse, neighboring quadrics and heap keys are updated. QEM is fast, produces high-quality results, and is the standard for offline LOD generation.

\textbf{Q3. How do you compute the quadric matrix Q for a vertex?}\\
\textbf{Steps:} 1) For each face $f$ adjacent to vertex $v$, compute its plane equation $ax+by+cz+d=0$ from the unit normal and one vertex. 2) Form the \textbf{fundamental quadric} $K_f = \mathbf{p}\mathbf{p}^T$, a symmetric $4\times 4$ matrix. 3) Sum across the one-ring: $Q_v = \sum_{f \in N(v)} K_f$. The diagonal/upper triangle has 10 unique entries, so storage is compact. Optionally weight $K_f$ by face \textbf{area} to give larger triangles more influence. For boundary edges, add a \textbf{virtual perpendicular plane} with high weight to prevent boundary shrinkage. The resulting $Q$ encodes the squared-distance field to the local tangent planes.

\textbf{Q4. How do you choose the optimal position v* after collapsing edge $(v_i,v_j)$?}\\
Minimize $\Delta(v)=v^T \bar Q v$ where $\bar Q = Q_i + Q_j$. Setting $\partial \Delta / \partial v = 0$ gives the linear system: solve
\[
\begin{pmatrix} q_{11} & q_{12} & q_{13} & q_{14}\\ q_{12} & q_{22} & q_{23} & q_{24}\\ q_{13} & q_{23} & q_{33} & q_{34}\\ 0 & 0 & 0 & 1 \end{pmatrix} v^* = \begin{pmatrix} 0\\0\\0\\1 \end{pmatrix}
\]
i.e.\ invert the upper-left $3\times 3$ block. If the matrix is \textbf{singular} (degenerate planes), fall back to: (a) midpoint $(v_i+v_j)/2$, (b) endpoint with smaller $\Delta$, or (c) line search along the edge. The optimal $v^*$ generally lies off the edge, allowing better fit to curved surfaces.

\textbf{Q5. What is progressive mesh? How does it enable streaming and geomorphing?}\\
A \textbf{progressive mesh} (Hoppe 1996) stores a base coarse mesh $M_0$ plus an ordered sequence of \textbf{vertex splits} (the inverse of edge collapses): $M_0 \to M_1 \to \ldots \to M_n$. Each split adds one vertex and two faces. \textbf{Streaming:} the client renders $M_0$ immediately while remaining splits arrive over the network, refining detail progressively. \textbf{Geomorphing:} during LOD switching, vertex positions are linearly interpolated between $M_k$ and $M_{k+1}$ over a few frames, eliminating the visual ``popping'' artifact. Progressive meshes also support \textbf{selective refinement}, refining only regions near the camera. They are widely used in games and 3D web (glTF Draco extensions adopt similar ideas).

\textbf{Q6. Compare vertex clustering vs vertex decimation vs QEM edge collapse.}\\
\textbf{Vertex clustering} (Rossignac-Borrel): overlay a 3D grid, collapse all vertices in each cell to one representative. Very fast, out-of-core friendly, but loses topology and quality. \textbf{Vertex decimation} (Schroeder): iteratively delete a vertex if its local neighborhood is ``flat'' (distance-to-plane threshold) and re-triangulate the hole. Preserves topology, decent quality, but greedy and not error-bounded. \textbf{QEM edge collapse} (Garland-Heckbert): each collapse minimizes $v^T Q v$, giving provably good error. Slower than clustering, faster than decimation per quality level, and is the de-facto standard. Quality: QEM > Decimation > Clustering. Speed: Clustering > QEM > Decimation.

\textbf{Q7. What metrics measure simplification quality?}\\
\textbf{Hausdorff distance} $d_H(A,B)=\max\{\sup_{a\in A}\inf_{b\in B}\|a-b\|, \sup_{b\in B}\inf_{a\in A}\|a-b\|\}$ — worst-case symmetric distance, very sensitive to outliers. \textbf{RMS error} $\sqrt{\tfrac{1}{|S|}\sum_{p\in S} d(p,M')^2}$ over sampled points $S$ on the original mesh — average-case, smoother metric. Tools like \textbf{Metro} and \textbf{MeshLab} compute both. Other measures: \textbf{normal deviation} (visual smoothness), \textbf{volume preservation}, and perceptual metrics (\textbf{MSDM2}). For texture-mapped models, also track UV/texture distortion. Report errors as a percentage of the bounding-box diagonal so they are scale-invariant.

\textbf{Q8. What is LOD (level of detail)? How does it relate to simplification?}\\
\textbf{LOD} renders different mesh resolutions depending on distance, screen size, or importance. \textbf{Discrete LOD} stores a few precomputed meshes (e.g., 100\%, 50\%, 25\%, 10\%) and switches based on distance thresholds. \textbf{Continuous LOD} (CLOD) uses a progressive mesh and refines vertex-by-vertex. \textbf{Hierarchical LOD} (HLOD) groups multiple objects into one simplified proxy for distant clusters (Unreal Engine 4/5). Simplification algorithms (mainly QEM) generate the LOD chain offline. LOD reduces \textbf{vertex shader} load, draw-call cost, and memory bandwidth. Modern engines (Unreal Nanite) virtualize geometry into a meshlet hierarchy, blending LOD techniques with cluster culling.

\textbf{Q9. How is QEM used to create LOD models for VR/games?}\\
\textbf{Pipeline:} 1) Author or scan the high-poly hero asset (e.g., 1M triangles). 2) Run QEM offline to generate LOD0 (full), LOD1 (50\%), LOD2 (20\%), LOD3 (5\%), LOD4 (1\%). 3) \textbf{Bake} normal maps from LOD0 onto LOD1+ to recover surface detail visually. 4) Constrain QEM to \textbf{preserve UV seams, sharp edges, and material boundaries} via boundary-plane penalties. 5) Export with progressive mesh metadata if streaming. 6) The engine selects LOD by projected screen size each frame. For \textbf{VR} (90+ Hz, two eyes), lower LODs are crucial, and \textbf{foveated} LOD pushes higher detail only at gaze. QEM's quality-per-budget makes it ideal for tight VR frame times.

\textbf{Q10. What is feature-preserving simplification? Why is it important?}\\
\textbf{Feature-preserving simplification} explicitly retains \textbf{sharp edges, corners, silhouettes, and textures} that carry semantic meaning. Without it, mechanical parts lose chamfers, characters lose facial creases, and CAD models lose dimensions. Techniques: 1) Detect features via \textbf{dihedral angle} threshold ($> 60°$ = sharp). 2) Add high-weight \textbf{constraint planes} along feature edges to the QEM matrix. 3) Forbid collapses that would destroy a feature edge. 4) Preserve \textbf{boundary} loops similarly. 5) Use \textbf{appearance-based} QEM that incorporates color/UV gradients (Garland-Heckbert 1998). Important because human perception is highly sensitive to silhouettes and creases; even small errors there are immediately visible.

\textbf{Q11. What is noise in 3D data and where does it come from?}\\
\textbf{Noise} is unwanted random/systematic deviation of measured points from the true surface. Sources: 1) \textbf{Sensor noise} — laser speckle, photon shot noise in ToF/structured-light cameras (Kinect, RealSense). 2) \textbf{Calibration error} — between cameras and laser. 3) \textbf{Registration error} — when fusing multiple scans (ICP residuals). 4) \textbf{Surface properties} — specular, transparent, dark surfaces give bad returns. 5) \textbf{Environmental} — ambient light, vibration, motion blur. 6) \textbf{Discretization} — voxel grid quantization. Noise is usually modeled as \textbf{Gaussian} along the surface normal, but real LiDAR also has \textbf{outliers} (multi-path, mixed pixels) and structured artifacts (banding).

\textbf{Q12. Given a noisy point cloud, how do you make it more regular? List algorithms.}\\
\textbf{Pipeline:} 1) \textbf{Outlier removal}: Statistical Outlier Removal (SOR), Radius Outlier Removal. 2) \textbf{Downsample}: Voxel-grid filter or Poisson-disk sampling for uniform density. 3) \textbf{Smooth in place}: \textbf{MLS} (Moving Least Squares) projects each point to a local polynomial fit. 4) \textbf{Resample/regularize}: \textbf{WLOP} (Weighted Locally Optimal Projection) redistributes points uniformly even on noisy input. 5) \textbf{Bilateral point-cloud filter} preserves features while smoothing. 6) \textbf{Normal estimation} via PCA, then orient consistently (minimum spanning tree). 7) Optionally \textbf{reconstruct mesh} (Poisson, BPA) and apply mesh denoising. 8) Modern: deep models like \textbf{PointCleanNet}, \textbf{Score-Based} denoising.

\textbf{Q13. Explain Laplacian smoothing — formula and why it causes shrinkage.}\\
For each vertex $v_i$, the \textbf{umbrella operator}/Laplacian is
\[
L(v_i) = \frac{1}{|N(i)|}\sum_{j\in N(i)} (v_j - v_i)
\]
The update is $v_i \leftarrow v_i + \lambda L(v_i)$ with $0 < \lambda < 1$. It diffuses geometry like the heat equation $\partial v/\partial t = \Delta v$. \textbf{Shrinkage:} $L$ acts as a low-pass filter that attenuates all non-zero spatial frequencies, so after many iterations a closed surface converges toward its centroid (a sphere shrinks to a point). It also \textbf{erodes sharp features} because creases are high-frequency content. Use \textbf{Taubin} or pin boundary vertices to mitigate.

\textbf{Q14. Compare uniform Laplacian vs cotangent Laplacian — which is better?}\\
\textbf{Uniform Laplacian} weights each neighbor equally: $L_u(v_i)=\tfrac{1}{|N(i)|}\sum (v_j-v_i)$. Simple, fast, but depends on triangulation: irregular meshes drift tangentially (not just normally). \textbf{Cotangent Laplacian}: $L_c(v_i)=\tfrac{1}{2A_i}\sum_j (\cot\alpha_{ij}+\cot\beta_{ij})(v_j-v_i)$, where $\alpha,\beta$ are the angles opposite edge $(i,j)$ and $A_i$ the Voronoi area. This is the proper discretization of the \textbf{Laplace-Beltrami} operator and approximates \textbf{mean curvature normal}: $\Delta v = -2H\mathbf{n}$. Cotangent is better for smoothing, parameterization, and physics, because it is intrinsic and triangulation-independent (when angles stay non-obtuse). Drawback: cotangent weights can be \textbf{negative} in obtuse triangles, requiring clamping or remeshing.

\textbf{Q15. Explain Taubin $\lambda|\mu$ smoothing — why two passes prevent shrinkage.}\\
\textbf{Taubin (1995):} alternate two updates per iteration:
\[
v_i^{(t+1/2)} = v_i^{(t)} + \lambda L(v_i^{(t)}), \quad v_i^{(t+1)} = v_i^{(t+1/2)} + \mu L(v_i^{(t+1/2)})
\]
with $0 < \lambda < -\mu$, e.g.\ $\lambda=0.33, \mu=-0.34$. The first (positive $\lambda$) shrinks; the second (negative $\mu$) inflates. The combined transfer function is $f(k)=(1-\lambda k)(1-\mu k)$ in the spectral domain of the Laplacian. The \textbf{pass-band frequency} $k_{PB}=\tfrac{1}{\lambda}+\tfrac{1}{\mu}$ satisfies $f(k_{PB})=1$ — frequencies below $k_{PB}$ are preserved, higher ones attenuated. Choosing $\lambda+\mu < 0$ gives a true low-pass filter without DC shrinkage. No matrix solve needed; it's still a local averaging.

\textbf{Q16. Explain bilateral mesh denoising. How does it preserve features?}\\
\textbf{Bilateral mesh denoising} (Fleishman 2003, Jones 2003) extends 2D bilateral filtering to meshes. For vertex $v_i$ with normal $n_i$, predict its new position by displacing along $n_i$:
\[
v_i' = v_i + n_i \cdot \frac{\sum_{j\in N(i)} w_c(\|v_j-v_i\|)\, w_s(\langle n_i, v_j-v_i\rangle)\,\langle n_i, v_j-v_i\rangle}{\sum w_c w_s}
\]
where $w_c(x)=\exp(-x^2/2\sigma_c^2)$ is the \textbf{spatial} (closeness) weight and $w_s(y)=\exp(-y^2/2\sigma_s^2)$ is the \textbf{range} (similarity along normal) weight. \textbf{Feature preservation:} across a sharp edge, neighbors lie far from $v_i$'s tangent plane, so $w_s$ becomes small and they don't pull $v_i$ across — features stay sharp while flat regions still smooth. Tune $\sigma_c$ to neighborhood size and $\sigma_s$ to noise amplitude.

\switchcolumn

\section*{QA2b. Câu hỏi lý thuyết — Đơn giản hóa \& Khử nhiễu}

\textbf{Câu 1. Đơn giản hóa lưới là gì? Vì sao cần?}\\
\textbf{Đơn giản hóa lưới} (mesh simplification) là quá trình giảm số đỉnh, cạnh, mặt của lưới tam giác mà vẫn giữ hình dáng trực quan và hình học gần nguyên gốc nhất có thể. Máy quét và CAD hiện đại tạo ra lưới vài triệu tam giác — quá nặng để render real-time, truyền mạng, hay mô phỏng vật lý. Đơn giản hóa tạo mô hình nhẹ phù hợp với \textbf{LOD (mức độ chi tiết)}, thiết bị mobile/VR, và 3D web. Nó cũng loại bỏ chi tiết vi mô không đóng góp vào nhận thức. Thuật toán phổ biến: \textbf{vertex clustering}, \textbf{vertex decimation} (Schroeder), và \textbf{QEM edge collapse} (Garland-Heckbert). Mục tiêu: tối thiểu sai số hình học với ngân sách tam giác cho trước.

\textbf{Câu 2. Quadric Error Metric (QEM) edge collapse — viết công thức.}\\
Với mỗi đỉnh $v$, \textbf{quadric} là $Q = \sum_{p} K_p$ trong đó mỗi mặt phẳng $p: ax+by+cz+d=0$ đóng góp $K_p = \mathbf{p}\mathbf{p}^T$ với $\mathbf{p}=(a,b,c,d)^T$. Sai số khi đặt $v=(x,y,z,1)^T$ là tổng bình phương khoảng cách đến các mặt phẳng đó: $\Delta(v) = v^T Q v$. Khi \textbf{gộp cạnh} $(v_i,v_j)\to\bar v$: quadric mới $\bar Q = Q_i + Q_j$, sai số $\Delta(\bar v)=\bar v^T \bar Q \bar v$. Đẩy các cạnh vào heap khóa bằng sai số tối thiểu này và gộp theo thứ tự. Sau mỗi lần gộp, cập nhật quadric và khóa heap của hàng xóm. QEM nhanh, chất lượng cao, là chuẩn de-facto cho LOD offline.

\textbf{Câu 3. Cách tính ma trận quadric Q cho một đỉnh?}\\
\textbf{Các bước:} 1) Với mỗi mặt $f$ kề đỉnh $v$, tính phương trình mặt phẳng $ax+by+cz+d=0$ từ pháp tuyến đơn vị và một đỉnh. 2) Lập \textbf{quadric cơ bản} $K_f=\mathbf{p}\mathbf{p}^T$, ma trận đối xứng $4\times 4$. 3) Cộng dồn quanh \textbf{one-ring}: $Q_v=\sum_{f\in N(v)} K_f$. Vì đối xứng nên chỉ cần lưu 10 phần tử riêng biệt. Tùy chọn: nhân thêm \textbf{diện tích} mặt $f$ để tam giác lớn ảnh hưởng nhiều hơn. Với cạnh biên, thêm \textbf{mặt phẳng vuông góc ảo} có trọng số lớn để chống co biên. Kết quả $Q$ mã hóa trường khoảng cách bình phương đến các mặt phẳng tiếp tuyến cục bộ.

\textbf{Câu 4. Chọn vị trí tối ưu $v^*$ sau khi gộp cạnh $(v_i,v_j)$ thế nào?}\\
Tối thiểu hóa $\Delta(v)=v^T \bar Q v$ với $\bar Q=Q_i+Q_j$. Đạo hàm bằng 0 dẫn đến hệ tuyến tính:
\[
\begin{pmatrix} q_{11}&q_{12}&q_{13}&q_{14}\\ q_{12}&q_{22}&q_{23}&q_{24}\\ q_{13}&q_{23}&q_{33}&q_{34}\\ 0&0&0&1\end{pmatrix} v^* = \begin{pmatrix}0\\0\\0\\1\end{pmatrix}
\]
tức là nghịch đảo khối $3\times 3$ trên-trái. Nếu ma trận \textbf{suy biến} (mặt phẳng thoái hóa), dự phòng bằng: (a) trung điểm $(v_i+v_j)/2$, (b) đỉnh có $\Delta$ nhỏ hơn, hoặc (c) line search dọc cạnh. $v^*$ tối ưu thường nằm ngoài cạnh, cho phép khớp tốt hơn với mặt cong.

\textbf{Câu 5. Progressive mesh là gì? Nó hỗ trợ streaming và geomorphing ra sao?}\\
\textbf{Progressive mesh} (Hoppe 1996) lưu lưới thô gốc $M_0$ kèm chuỗi \textbf{vertex split} có thứ tự (đảo của edge collapse): $M_0\to M_1\to\ldots\to M_n$. Mỗi split thêm 1 đỉnh và 2 mặt. \textbf{Streaming:} client render $M_0$ ngay, các split còn lại đến qua mạng để tinh chỉnh dần. \textbf{Geomorphing:} khi chuyển LOD, nội suy tuyến tính vị trí đỉnh giữa $M_k$ và $M_{k+1}$ trong vài frame, tránh hiện tượng ``popping''. Progressive mesh còn hỗ trợ \textbf{selective refinement} — chỉ tinh chỉnh vùng gần camera. Được dùng rộng rãi trong games và 3D web (glTF Draco có ý tưởng tương tự).

\textbf{Câu 6. So sánh vertex clustering, vertex decimation, QEM edge collapse.}\\
\textbf{Vertex clustering} (Rossignac-Borrel): phủ lưới 3D, gộp tất cả đỉnh trong mỗi ô về một đại diện. Rất nhanh, chạy out-of-core được, nhưng mất topology và chất lượng kém. \textbf{Vertex decimation} (Schroeder): lặp lại xóa đỉnh nếu vùng quanh nó ``phẳng'' (ngưỡng khoảng cách-mặt phẳng) rồi tam giác hóa lại lỗ. Giữ topology, chất lượng khá, nhưng tham lam và không ràng buộc sai số. \textbf{QEM edge collapse} (Garland-Heckbert): mỗi lần gộp tối thiểu hóa $v^T Q v$, chứng minh được tốt. Chậm hơn clustering, nhanh hơn decimation cùng mức chất lượng. Chất lượng: QEM > Decimation > Clustering. Tốc độ: Clustering > QEM > Decimation.

\textbf{Câu 7. Các metric đo chất lượng đơn giản hóa?}\\
\textbf{Hausdorff distance} $d_H(A,B)=\max\{\sup_a\inf_b\|a-b\|, \sup_b\inf_a\|a-b\|\}$ — khoảng cách đối xứng tệ-nhất, nhạy với outlier. \textbf{RMS error} $\sqrt{\tfrac{1}{|S|}\sum d(p,M')^2}$ trên các điểm sampled $S$ trên lưới gốc — trung bình, mượt hơn. Công cụ \textbf{Metro} và \textbf{MeshLab} tính cả hai. Các đại lượng khác: \textbf{độ lệch pháp tuyến} (mượt thị giác), \textbf{bảo toàn thể tích}, và metric tri giác (\textbf{MSDM2}). Với mô hình có texture, thêm \textbf{distortion UV/texture}. Báo cáo sai số dưới dạng phần trăm đường chéo bounding-box để bất biến tỷ lệ.

\textbf{Câu 8. LOD là gì? Liên hệ với simplification?}\\
\textbf{LOD (level of detail)} render các độ phân giải khác nhau tùy khoảng cách, kích thước trên màn hình, hoặc mức quan trọng. \textbf{LOD rời rạc}: lưu vài lưới đã tính trước (100\%, 50\%, 25\%, 10\%) và chuyển theo ngưỡng khoảng cách. \textbf{LOD liên tục} (CLOD): dùng progressive mesh, tinh chỉnh từng đỉnh. \textbf{HLOD} gộp nhiều object thành proxy đơn giản cho cụm xa (Unreal Engine 4/5). Thuật toán simplification (chủ yếu QEM) sinh chuỗi LOD offline. LOD giảm tải \textbf{vertex shader}, draw call, băng thông bộ nhớ. Engine hiện đại (Unreal Nanite) ảo hóa hình học thành phân cấp meshlet, kết hợp LOD với cluster culling.

\textbf{Câu 9. QEM tạo LOD cho VR/games như thế nào?}\\
\textbf{Pipeline:} 1) Tạo asset hero high-poly (vd 1M tam giác). 2) Chạy QEM offline sinh LOD0 (đầy đủ), LOD1 (50\%), LOD2 (20\%), LOD3 (5\%), LOD4 (1\%). 3) \textbf{Bake} normal map từ LOD0 lên LOD1+ để giữ chi tiết bề mặt. 4) Ràng buộc QEM \textbf{giữ UV seam, cạnh sắc, biên material} bằng penalty mặt biên. 5) Export kèm metadata progressive nếu streaming. 6) Engine chọn LOD theo kích thước chiếu trên màn hình mỗi frame. Với \textbf{VR} (90+ Hz, 2 mắt), LOD thấp là then chốt; \textbf{foveated LOD} chỉ giữ chi tiết cao tại điểm nhìn. Tỷ lệ chất lượng/ngân sách của QEM rất phù hợp với frame time chặt của VR.

\textbf{Câu 10. Simplification giữ đặc trưng (feature-preserving) là gì? Tại sao quan trọng?}\\
\textbf{Simplification giữ đặc trưng} cố tình giữ lại \textbf{cạnh sắc, góc, đường bao silhouette, texture} — những thứ mang ý nghĩa ngữ nghĩa. Không có nó, chi tiết cơ khí mất chamfer, nhân vật mất nếp mặt, mô hình CAD mất kích thước. Kỹ thuật: 1) Phát hiện feature qua ngưỡng \textbf{góc nhị diện} ($>60°$ = sắc). 2) Thêm \textbf{mặt phẳng ràng buộc} trọng số cao dọc cạnh feature vào ma trận QEM. 3) Cấm collapse phá hủy cạnh feature. 4) Bảo toàn vòng \textbf{biên} tương tự. 5) Dùng QEM \textbf{appearance-based} tích hợp gradient màu/UV (Garland-Heckbert 1998). Quan trọng vì thị giác con người rất nhạy với silhouette và crease — sai số nhỏ ở đó cũng dễ thấy.

\textbf{Câu 11. Nhiễu trong dữ liệu 3D là gì, đến từ đâu?}\\
\textbf{Nhiễu} là độ lệch ngẫu nhiên/hệ thống của điểm đo so với bề mặt thật. Nguồn: 1) \textbf{Sensor noise} — speckle laser, shot noise photon trong camera ToF/structured-light (Kinect, RealSense). 2) \textbf{Calibration} — sai lệch giữa camera và laser. 3) \textbf{Đăng ký} — residual khi ICP gộp nhiều scan. 4) \textbf{Tính chất bề mặt} — bóng, trong suốt, tối cho tín hiệu kém. 5) \textbf{Môi trường} — ánh sáng, rung, motion blur. 6) \textbf{Lượng tử hóa} voxel grid. Nhiễu thường mô hình hóa \textbf{Gauss} dọc pháp tuyến, nhưng LiDAR thực còn có \textbf{outlier} (multi-path, mixed pixel) và artifact có cấu trúc (banding).

\textbf{Câu 12. Có một point cloud nhiễu, làm sao chuẩn hóa? Liệt kê thuật toán.}\\
\textbf{Pipeline:} 1) \textbf{Loại outlier}: Statistical Outlier Removal (SOR), Radius Outlier Removal. 2) \textbf{Downsample}: Voxel-grid hoặc Poisson-disk để mật độ đều. 3) \textbf{Làm mượt tại chỗ}: \textbf{MLS} chiếu mỗi điểm lên đa thức cục bộ. 4) \textbf{Resample}: \textbf{WLOP} phân bố lại điểm đều ngay cả trên dữ liệu nhiễu. 5) \textbf{Bilateral} cho point cloud, giữ feature khi khử nhiễu. 6) \textbf{Ước lượng pháp tuyến} bằng PCA, sau đó định hướng nhất quán (cây bao trùm tối thiểu). 7) Tùy chọn \textbf{tái dựng lưới} (Poisson, BPA) rồi áp dụng denoising lưới. 8) Hiện đại: deep learning như \textbf{PointCleanNet}, \textbf{Score-Based} denoising.

\textbf{Câu 13. Laplacian smoothing — công thức và vì sao gây co?}\\
Với mỗi đỉnh $v_i$, toán tử \textbf{umbrella}/Laplacian là
\[
L(v_i)=\frac{1}{|N(i)|}\sum_{j\in N(i)}(v_j-v_i)
\]
Cập nhật: $v_i\leftarrow v_i+\lambda L(v_i)$ với $0<\lambda<1$. Nó khuếch tán hình học như phương trình nhiệt $\partial v/\partial t=\Delta v$. \textbf{Co (shrinkage):} $L$ là bộ lọc thông thấp triệt tiêu mọi tần số không gian khác 0, nên sau nhiều lần lặp, mặt kín hội tụ về tâm khối lượng (hình cầu co thành điểm). Nó cũng \textbf{xói mòn feature} vì crease là tần số cao. Khắc phục: dùng \textbf{Taubin} hoặc ghim đỉnh biên.

\textbf{Câu 14. Uniform Laplacian vs cotangent Laplacian — cái nào tốt hơn?}\\
\textbf{Uniform}: trọng số bằng nhau $L_u(v_i)=\tfrac{1}{|N(i)|}\sum(v_j-v_i)$. Đơn giản, nhanh, nhưng phụ thuộc tam giác hóa: lưới không đều bị trôi tiếp tuyến (không chỉ pháp tuyến). \textbf{Cotangent}: $L_c(v_i)=\tfrac{1}{2A_i}\sum_j(\cot\alpha_{ij}+\cot\beta_{ij})(v_j-v_i)$, $\alpha,\beta$ là góc đối cạnh $(i,j)$, $A_i$ là diện tích Voronoi. Đây là rời rạc hóa đúng của \textbf{Laplace-Beltrami}, xấp xỉ \textbf{vector pháp tuyến nhân độ cong trung bình}: $\Delta v=-2H\mathbf{n}$. Cotangent tốt hơn cho làm mượt, parameterization, vật lý — vì nội tại, không phụ thuộc tam giác hóa (khi góc không tù). Nhược: trọng số cotangent có thể \textbf{âm} ở tam giác tù, cần clamp hoặc remesh.

\textbf{Câu 15. Taubin $\lambda|\mu$ — vì sao 2 pass chống co? Điều kiện pass-band.}\\
\textbf{Taubin (1995):} mỗi vòng lặp xen kẽ 2 cập nhật:
\[
v^{(t+1/2)}=v^{(t)}+\lambda L(v^{(t)}),\quad v^{(t+1)}=v^{(t+1/2)}+\mu L(v^{(t+1/2)})
\]
với $0<\lambda<-\mu$, vd $\lambda=0.33,\mu=-0.34$. Pass đầu ($\lambda$ dương) co; pass sau ($\mu$ âm) phồng. Hàm truyền tổng hợp trong miền phổ Laplacian là $f(k)=(1-\lambda k)(1-\mu k)$. \textbf{Tần số pass-band} $k_{PB}=\tfrac{1}{\lambda}+\tfrac{1}{\mu}$ thỏa $f(k_{PB})=1$ — tần số nhỏ hơn $k_{PB}$ được giữ, lớn hơn bị suy giảm. Chọn $\lambda+\mu<0$ cho bộ lọc thông thấp thực sự, không co DC. Không cần giải hệ ma trận; vẫn chỉ là trung bình cục bộ.

\textbf{Câu 16. Bilateral mesh denoising. Tại sao giữ được feature?}\\
\textbf{Bilateral mesh denoising} (Fleishman 2003, Jones 2003) mở rộng bilateral 2D sang lưới. Với đỉnh $v_i$ pháp tuyến $n_i$, dự đoán vị trí mới bằng dịch dọc $n_i$:
\[
v_i'=v_i+n_i\cdot\frac{\sum_j w_c(\|v_j-v_i\|)\,w_s(\langle n_i,v_j-v_i\rangle)\,\langle n_i,v_j-v_i\rangle}{\sum w_c w_s}
\]
$w_c(x)=\exp(-x^2/2\sigma_c^2)$ là trọng số \textbf{khoảng cách}, $w_s(y)=\exp(-y^2/2\sigma_s^2)$ là trọng số \textbf{độ lệch dọc pháp tuyến}. \textbf{Giữ feature:} qua cạnh sắc, hàng xóm cách xa mặt phẳng tiếp tuyến của $v_i$, $w_s$ bé, không kéo $v_i$ qua bên kia — feature giữ nguyên còn vùng phẳng vẫn được làm mượt. Tinh chỉnh $\sigma_c$ theo kích thước neighborhood, $\sigma_s$ theo biên độ nhiễu.

\switchcolumn*

\textbf{Q17. What is mean curvature flow smoothing?}\\
\textbf{Mean curvature flow} (MCF) evolves the surface along its normal at a speed proportional to mean curvature: $\partial v/\partial t = -H\mathbf{n}$. Discretely, this equals applying the \textbf{cotangent Laplacian}: $\Delta v = -2H\mathbf{n}$, so MCF is $v_{t+1}=v_t+\tau\cdot \tfrac{1}{2}L_c(v_t)$. It smooths surfaces optimally in a geometric sense (curvature-shortening). \textbf{Implicit MCF} (Desbrun 1999) solves $(I-\tau L_c)v_{t+1}=v_t$ for unconditional stability with large $\tau$. Mean curvature flow is intrinsic — purely geometric, not topology-dependent. Drawback: still shrinks closed surfaces (sphere collapses to point); use \textbf{volume-preserving} or \textbf{conformalized} MCF (Kazhdan 2012) to keep size.

\textbf{Q18. Compare Laplacian vs Taubin vs bilateral mesh denoising.}\\
\textbf{Laplacian (1 pass):} simplest, fast, \textit{shrinks} the model and \textit{blurs features}. Use only on small noise. \textbf{Taubin $\lambda|\mu$:} no shrinkage (low-pass with proper pass-band), still feature-blurring because it is linear. Fast, no matrix solve. \textbf{Bilateral:} non-linear, weights neighbors by both spatial and normal-difference Gaussians — \textit{feature preserving} (sharp creases survive). Slower and needs tuning $\sigma_c, \sigma_s$. Modern: \textbf{normal-based bilateral} (Zheng 2011) filters face normals first then updates vertex positions; even better. Rule of thumb: smooth/organic data $\to$ Taubin; CAD/feature-rich $\to$ bilateral or $\ell_0$ minimization.

\textbf{Q19. Explain Statistical Outlier Removal (SOR) for point clouds.}\\
\textbf{SOR} (PCL implementation) removes points whose mean distance to their $k$ nearest neighbors deviates too far from the global distribution. \textbf{Steps:} 1) For each point $p_i$, find its $k$ nearest neighbors and compute the mean distance $d_i$. 2) Compute the global mean $\mu$ and standard deviation $\sigma$ of $\{d_i\}$. 3) A point is an \textbf{outlier} if $d_i > \mu + \alpha\sigma$ (typically $\alpha=1$ to $2$). 4) Remove all outliers. SOR assumes the overall density is roughly Gaussian and works well for sparse outliers from sensor noise. Parameters: $k$ controls locality (e.g., 30), $\alpha$ controls aggressiveness. It does \textit{not} preserve fine detail equally — dense/sparse boundaries can be misclassified.

\textbf{Q20. What is Radius Outlier Removal? When to use vs SOR?}\\
\textbf{Radius Outlier Removal} (ROR) keeps only points that have at least $N_{\min}$ neighbors within radius $r$. \textbf{Steps:} 1) For each point, range-search neighbors within $r$. 2) If count $< N_{\min}$, remove. Simpler than SOR, no statistics needed. \textbf{When to use:} 1) When you know the expected sensor density (e.g., LiDAR at 10 m gives $\sim$1 cm spacing). 2) For very sparse outliers (single floaters). 3) Real-time pipelines that can't afford SOR's two-pass statistics. \textbf{SOR vs ROR:} SOR adapts to local density, ROR uses fixed thresholds. Use SOR when density varies (close vs far objects); use ROR when density is uniform and speed matters.

\textbf{Q21. What is voxel grid downsampling?}\\
\textbf{Voxel grid downsampling} overlays a regular 3D grid of cubes (voxels) of size $L$ on the point cloud and replaces all points in each voxel with their \textbf{centroid} (or one representative). \textbf{Steps:} 1) Compute bounding box, divide by leaf size. 2) Hash each point into its voxel. 3) Average points (and normals/colors) per voxel. Output: roughly uniform-density cloud, fewer points. Benefits: fast, cache-friendly, reduces memory for downstream ICP/registration. Drawbacks: fixed grid causes \textbf{aliasing} and shifts at voxel boundaries. Variants: \textbf{octree-based} adaptive voxels, or \textbf{Poisson-disk} sampling for blue-noise distribution. Used in PCL (\verb|pcl::VoxelGrid|), Open3D, and ROS perception stacks.

\textbf{Q22. What is MLS (Moving Least Squares)? How is it used for denoising?}\\
\textbf{Moving Least Squares} (Levin 1998, Alexa 2003) projects each input point to a smooth surface defined locally. \textbf{Steps:} 1) For point $p$, gather neighbors within radius $h$. 2) Fit a \textbf{local reference plane} by weighted PCA with Gaussian weights $w(\|q-p\|)=\exp(-\|q-p\|^2/h^2)$. 3) On that plane, fit a low-degree \textbf{polynomial} (often quadratic) approximating the height field of neighbors. 4) Project $p$ onto this polynomial to get $p'$. The MLS surface is smooth (infinitely differentiable in $h$) and noise is averaged out at scale $h$. Used in PCL (\verb|pcl::MovingLeastSquares|), Meshlab. Drawbacks: bandwidth $h$ trade-off — too small keeps noise, too large blurs features.

\textbf{Q23. What is WLOP? How does it differ from MLS?}\\
\textbf{WLOP} (Weighted Locally Optimal Projection, Huang 2009) generates a uniformly distributed, denoised set of \textbf{output} points $\{x_i\}$ from noisy input $\{p_j\}$, even \textit{without} normal estimation. It minimizes:
\[
E(X) = \sum_i \sum_j \theta(\|x_i-p_j\|)\|x_i-p_j\| - \mu \sum_i \sum_{i'\neq i} \theta(\|x_i-x_{i'}\|)\eta(\|x_i-x_{i'}\|)
\]
The first \textbf{attraction} term pulls $x_i$ toward input data. The second \textbf{repulsion} term spreads output points apart for uniform density. Local density weights $\theta$ down-weight crowded regions. \textbf{Differences from MLS:} (a) MLS projects each input point individually and may inherit non-uniformity; WLOP redistributes; (b) WLOP needs no normals; (c) MLS needs a smooth surface assumption, WLOP handles outlier-rich data better. Use WLOP for raw scanner output that is highly non-uniform.

\textbf{Q24. How do you estimate normals for a point cloud (PCA-based)?}\\
\textbf{Steps:} 1) For point $p_i$, find $k$-NN or radius-$r$ neighbors $\{q_j\}$. 2) Compute the centroid $\bar p$ and \textbf{covariance matrix} $C=\tfrac{1}{k}\sum (q_j-\bar p)(q_j-\bar p)^T$. 3) Eigen-decompose: $C=V\Lambda V^T$ with $\lambda_0\le\lambda_1\le\lambda_2$. The eigenvector for the smallest $\lambda_0$ is the normal $n_i$ (perpendicular to the local fitted plane). 4) The ratio $\lambda_0/(\lambda_0+\lambda_1+\lambda_2)$ is the \textbf{surface variation}/curvature proxy. 5) \textbf{Orient consistently}: compute a Riemannian graph of nearest neighbors, build a minimum spanning tree, and propagate orientation by flipping if $n_i\cdot n_j<0$ (Hoppe 1992). Robust variants weight neighbors by distance or use RANSAC plane fitting.

\textbf{Q25. Compare mesh smoothing vs point cloud smoothing — what's different?}\\
\textbf{Mesh smoothing} relies on explicit \textbf{connectivity}: each vertex has a fixed one-ring of neighbors, so Laplacian/Taubin/cotangent operators are well-defined and feature-preserving filters can use dihedral angles. \textbf{Point cloud smoothing} has no connectivity — you must reconstruct neighborhoods on-the-fly via $k$-NN or radius search every time. Operations: MLS, WLOP, bilateral PC, or learned models. Mesh smoothing also handles \textbf{topology} and can guarantee manifoldness; PC methods don't preserve any topology. Mesh has explicit normals from face geometry; PC needs PCA estimation that itself can be noisy. PC smoothing must additionally manage \textbf{density variation} (close vs far). Often the pipeline is: smooth PC $\to$ reconstruct mesh $\to$ smooth mesh.

\textbf{Q26. Describe the geometric modeling pipeline end-to-end.}\\
\textbf{Stages:} 1) \textbf{Acquisition} — laser scan, photogrammetry, depth camera, or CAD authoring. 2) \textbf{Registration} — align multiple scans (ICP, global SLAM). 3) \textbf{Cleaning} — remove outliers (SOR/ROR), fill gaps. 4) \textbf{Downsample} — voxel grid, Poisson-disk. 5) \textbf{Normal estimation} — PCA + consistent orientation. 6) \textbf{Surface reconstruction} — Poisson, Ball-Pivoting, Marching Cubes (from implicit). 7) \textbf{Mesh repair} — fill holes, fix non-manifold edges. 8) \textbf{Denoising} — Taubin, bilateral, MCF. 9) \textbf{Simplification \& LOD} — QEM, progressive mesh. 10) \textbf{Parameterization \& UV} — for texturing. 11) \textbf{Export} — OBJ/PLY/STL/glTF. 12) \textbf{Rendering / Application} — game engine, 3D printer, FEA, VR.

\textbf{Q27. Common 3D file formats — compare OBJ, PLY, STL, OFF.}\\
\textbf{OBJ} (Wavefront): ASCII, supports vertices, normals, UVs, materials (\verb|.mtl|), polygon faces (not just triangles). De-facto interchange for art assets. \textbf{PLY} (Stanford): ASCII or binary, supports arbitrary per-vertex/per-face properties (color, intensity, custom). Standard for scanner output and research. \textbf{STL}: ASCII or binary, only triangle soup with face normals — \textit{no shared vertex indexing, no color, no UV}. The standard for 3D printing/CAD slicers. \textbf{OFF}: ASCII, very simple — header line, vertex list, face list with index. Used in academia (Princeton Shape Benchmark). Also worth knowing: \textbf{glTF 2.0} (binary, PBR, animations — modern web/games), \textbf{FBX} (proprietary, animation-rich), \textbf{USD} (Pixar, scene graph).

\textbf{Q28. What are the main 3D representations?}\\
1) \textbf{Point cloud} — unordered set of $(x,y,z)$ points; native for scanners; no topology. 2) \textbf{Triangle/poly mesh} — vertices + edges + faces; explicit connectivity; standard for rendering. 3) \textbf{NURBS} — non-uniform rational B-spline surfaces; smooth, parametric, exact for CAD. 4) \textbf{Implicit / SDF} — $f(x,y,z)=0$; supports easy CSG, constant-time inside test; needs Marching Cubes to mesh. 5) \textbf{Voxel grid} — 3D pixels; native for medical/volumetric data; high memory. 6) \textbf{CSG} — Constructive Solid Geometry, boolean tree of primitives; classic CAD. 7) \textbf{Subdivision surfaces} — coarse mesh + refinement rules; smooth at any LOD. 8) \textbf{Neural fields} (NeRF, occupancy networks) — MLP encodes geometry/radiance.

\textbf{Q29. When would you choose mesh vs implicit vs voxel representation?}\\
\textbf{Mesh:} best for \textit{rendering} and \textit{games} — GPU loves triangles, surface-only memory, explicit UV. Bad for boolean ops and topology change. \textbf{Implicit / SDF:} best for \textit{modeling}, \textit{CSG booleans}, \textit{ray-marching} (shadertoy), \textit{collision detection} — closed-form inside test. Bad for direct rendering without conversion. \textbf{Voxel:} best for \textit{volumetric data} (medical CT/MRI), \textit{simulation} (fluids, voxel-based physics), \textit{games like Minecraft}. Bad for memory: $O(n^3)$ growth — solved by sparse octrees (\textbf{VDB}, \textbf{NanoVDB}). \textbf{Rule of thumb:} render = mesh, sculpt = SDF, simulate = voxel. Modern engines mix: Unreal Nanite uses mesh, Lumen uses SDF, fluid sims use VDB.

\textbf{Q30. What is $C^0, C^1, C^2$ continuity? Why does it matter?}\\
At the boundary between two curves/surfaces: \textbf{$C^0$} = positions match (no gap). \textbf{$C^1$} = positions and first derivatives (\textbf{tangents}) match — surfaces visually smooth, no ``kink''. \textbf{$C^2$} = positions, tangents, and second derivatives (\textbf{curvature}) match — even smoother, important for highlights and reflections. \textbf{$G^n$} (geometric continuity) is the looser version: tangent/curvature \textit{directions} match, magnitudes can differ. \textbf{Why it matters:} 1) Class-A surfaces in automotive/product design need $G^2$ for clean reflections. 2) Animation paths need at least $C^1$ to avoid jerk. 3) Subdivision schemes (Catmull-Clark) produce $C^2$ except at extraordinary vertices. 4) Bezier patches need careful boundary tangent matching for $C^1$.

\textbf{Q31. How is point cloud denoising used in LiDAR / autonomous driving?}\\
LiDAR returns are noisy and contain spurious points from \textbf{rain, fog, snow, exhaust, dust, sensor multipath}. Pipeline in AVs (Waymo, Tesla, Mobileye): 1) \textbf{Range filter} — remove very near/far returns. 2) \textbf{ROR/SOR} — remove single floating points (rain droplets). 3) \textbf{Ground plane fit} (RANSAC) — separate ground from objects. 4) \textbf{Voxel downsample} — for real-time 3D detectors (PointPillars, CenterPoint). 5) \textbf{Temporal denoising} — accumulate frames, vote on persistence (random rain points won't reappear). 6) \textbf{Weather-specific filters} — DROR (Dynamic Radius Outlier Removal) for snow. Clean point clouds are critical for downstream \textbf{3D object detection}, \textbf{tracking}, and \textbf{SLAM}. Even small false positives can trigger phantom braking.

\textbf{Q32. How is mesh simplification used in real-time rendering / games?}\\
Games can't render 100M-triangle assets at 60-120 Hz. \textbf{Uses of simplification:} 1) \textbf{Discrete LOD chains} (LOD0-4) generated by QEM offline, switched by screen size. 2) \textbf{HLOD} — distant clusters of buildings merged into single proxy meshes (Unreal). 3) \textbf{Imposters} — even simpler than meshes, just billboards from extreme LOD. 4) \textbf{Shadow LOD} — separate, more aggressive simplification for shadow casters. 5) \textbf{Collision LOD} — simplified physics meshes for runtime. 6) \textbf{Streaming} — progressive mesh / Nanite virtualized geometry. 7) \textbf{VR/mobile} — aggressive LOD because thermal/perf budget is tight. The combination saves \textbf{vertex shader cost}, draw calls, memory bandwidth, and shadow-pass cost. Without simplification, modern open-world games would be impossible.

\textbf{Q33. How is denoising used in medical 3D scanning?}\\
Medical CT/MRI/ultrasound and surface scanners (intra-oral dental, body scanners) produce noisy volumes/meshes. \textbf{Uses:} 1) \textbf{CT reconstruction} — iterative algorithms (SART, MBIR) include denoising priors (TV, BM3D). 2) \textbf{Surface extraction} — Marching Cubes followed by Taubin smoothing for organ surfaces. 3) \textbf{Dental crown design} — bilateral mesh denoising preserves cusp/groove features while removing scanner speckle. 4) \textbf{Surgical planning} — point cloud denoising on registered intra-op scans for AR overlay. 5) \textbf{Implant/prosthesis fitting} — clean meshes for 3D printing. 6) \textbf{Radiomics / AI} — denoised images give more reliable features for diagnosis. Denoising must be \textbf{feature-preserving} — it cannot blur tumor edges or vessel walls, which is why bilateral and edge-aware methods dominate.

\textbf{Q34. Why is geometric modeling fundamental to digital twins?}\\
A \textbf{digital twin} is a synchronized 3D virtual replica of a physical asset (factory, building, city, body). Geometric modeling provides: 1) \textbf{Acquisition} — laser scan / photogrammetry to capture as-built reality. 2) \textbf{Cleaning + denoising} — usable model from raw scans. 3) \textbf{Simplification + LOD} — twins of cities have billions of triangles; LOD enables interactive use. 4) \textbf{Mesh + parameterization} — for sensor placement, simulations (CFD, FEM), thermal analysis. 5) \textbf{Continuous update} — re-scan, register, diff with previous mesh to detect change (e.g.\ structural deformation). 6) \textbf{Interoperability} — standard formats (USD, IFC, glTF) for sharing across tools. Without geometric modeling, digital twins would just be CAD drawings — they'd lack the high-fidelity, sensor-driven, time-evolving nature that makes them ``twins''.

\textbf{Q35. How is QEM combined with Marching Cubes in medical visualization?}\\
\textbf{Pipeline:} 1) Acquire CT/MRI volume. 2) Segment the organ via thresholding/AI to get binary or scalar field. 3) Run \textbf{Marching Cubes} (Lorensen 1987) to extract iso-surface mesh — typically \textit{millions of small triangles}, often staircase-aliased. 4) Apply \textbf{Taubin smoothing} to remove voxel staircase artifacts. 5) Run \textbf{QEM simplification} to reduce to tens of thousands of triangles, preserving anatomical features (vessel curvature, tumor shape). 6) Compute \textbf{normals + UV} for shading. 7) Render with PBR for surgical planning, VR/AR overlay, or 3D print. The combo is essential because Marching Cubes outputs are too dense for interactive VR but contain the correct topology — QEM + smoothing make them practical while QEM with feature-preservation keeps clinically critical detail (e.g., vessel branching, tumor margins).

\switchcolumn

\textbf{Câu 17. Mean curvature flow smoothing là gì?}\\
\textbf{Mean curvature flow} (MCF) tiến hóa bề mặt dọc pháp tuyến với tốc độ tỷ lệ độ cong trung bình: $\partial v/\partial t=-H\mathbf{n}$. Rời rạc, đây chính là áp dụng \textbf{cotangent Laplacian}: $\Delta v=-2H\mathbf{n}$, nên MCF là $v_{t+1}=v_t+\tau\cdot\tfrac{1}{2}L_c(v_t)$. Làm mượt tối ưu theo nghĩa hình học (curvature-shortening). \textbf{Implicit MCF} (Desbrun 1999) giải $(I-\tau L_c)v_{t+1}=v_t$ để ổn định vô điều kiện với $\tau$ lớn. MCF là nội tại — thuần hình học, không phụ thuộc topology. Nhược: vẫn co mặt kín (cầu sụp về điểm); dùng MCF \textbf{bảo toàn thể tích} hoặc \textbf{conformalized} (Kazhdan 2012) để giữ kích thước.

\textbf{Câu 18. So sánh Laplacian vs Taubin vs bilateral mesh denoising.}\\
\textbf{Laplacian (1 pass):} đơn giản, nhanh, \textit{co} mô hình và \textit{làm mờ feature}. Chỉ dùng khi nhiễu nhỏ. \textbf{Taubin $\lambda|\mu$:} không co (thông thấp với pass-band đúng), vẫn làm mờ feature vì tuyến tính. Nhanh, không cần giải ma trận. \textbf{Bilateral:} phi tuyến, có trọng số cả khoảng cách lẫn độ lệch pháp tuyến — \textit{giữ feature} (cạnh sắc tồn tại). Chậm hơn, cần tinh chỉnh $\sigma_c,\sigma_s$. Hiện đại: \textbf{normal-based bilateral} (Zheng 2011) lọc pháp tuyến mặt trước rồi cập nhật vị trí; còn tốt hơn. Nguyên tắc: dữ liệu hữu cơ/mượt $\to$ Taubin; CAD/nhiều feature $\to$ bilateral hoặc tối thiểu hóa $\ell_0$.

\textbf{Câu 19. Statistical Outlier Removal (SOR) cho point cloud?}\\
\textbf{SOR} (PCL) loại bỏ điểm có khoảng cách trung bình đến $k$ hàng xóm gần nhất lệch quá xa phân phối toàn cục. \textbf{Các bước:} 1) Với mỗi điểm $p_i$, tìm $k$-NN và tính khoảng cách trung bình $d_i$. 2) Tính trung bình toàn cục $\mu$ và độ lệch chuẩn $\sigma$ của $\{d_i\}$. 3) Điểm là \textbf{outlier} nếu $d_i>\mu+\alpha\sigma$ (thường $\alpha=1$ đến $2$). 4) Xóa outlier. SOR giả định mật độ tổng thể xấp xỉ Gauss và tốt với outlier thưa do nhiễu sensor. Tham số: $k$ kiểm soát locality (vd 30), $\alpha$ kiểm soát độ aggressive. SOR \textit{không} giữ chi tiết nhỏ tốt như nhau — biên dày/thưa có thể bị phân loại sai.

\textbf{Câu 20. Radius Outlier Removal là gì? Khi nào dùng so với SOR?}\\
\textbf{Radius Outlier Removal} (ROR) chỉ giữ điểm có ít nhất $N_{\min}$ hàng xóm trong bán kính $r$. \textbf{Các bước:} 1) Range-search neighbors trong $r$ cho mỗi điểm. 2) Nếu count $<N_{\min}$, xóa. Đơn giản hơn SOR, không cần thống kê. \textbf{Khi nào dùng:} 1) Khi biết mật độ sensor kỳ vọng (vd LiDAR ở 10m cho khoảng cách $\sim$1cm). 2) Cho outlier rất thưa (điểm trôi đơn lẻ). 3) Pipeline real-time không đủ thời gian SOR 2-pass. \textbf{SOR vs ROR:} SOR thích nghi với mật độ cục bộ, ROR dùng ngưỡng cố định. Dùng SOR khi mật độ biến thiên (vật gần vs xa); dùng ROR khi mật độ đều và cần nhanh.

\textbf{Câu 21. Voxel grid downsampling là gì?}\\
\textbf{Voxel grid downsampling} phủ lưới khối lập phương đều (voxel) cạnh $L$ lên đám mây điểm và thay tất cả điểm trong mỗi voxel bằng \textbf{centroid} (hoặc đại diện). \textbf{Các bước:} 1) Tính bounding box, chia theo leaf size. 2) Hash mỗi điểm vào voxel của nó. 3) Trung bình điểm (và pháp tuyến/màu) mỗi voxel. Output: cloud mật độ xấp xỉ đều, ít điểm hơn. Lợi: nhanh, cache-friendly, giảm bộ nhớ cho ICP/registration sau đó. Hại: lưới cố định gây \textbf{aliasing} và dịch ở biên voxel. Biến thể: \textbf{voxel thích nghi} dựa trên octree, hoặc \textbf{Poisson-disk} cho phân bố blue-noise. Dùng trong PCL (\verb|pcl::VoxelGrid|), Open3D, ROS perception stacks.

\textbf{Câu 22. MLS là gì? Dùng để khử nhiễu thế nào?}\\
\textbf{Moving Least Squares} (Levin 1998, Alexa 2003) chiếu mỗi điểm input về một bề mặt mượt định nghĩa cục bộ. \textbf{Các bước:} 1) Với điểm $p$, lấy hàng xóm trong bán kính $h$. 2) Khớp \textbf{mặt phẳng tham chiếu} cục bộ bằng PCA có trọng số $w(\|q-p\|)=\exp(-\|q-p\|^2/h^2)$. 3) Trên mặt phẳng đó, khớp \textbf{đa thức} bậc thấp (thường bậc 2) xấp xỉ trường độ cao của hàng xóm. 4) Chiếu $p$ lên đa thức để được $p'$. Mặt MLS mượt (khả vi vô hạn theo $h$) và nhiễu được trung bình hóa ở scale $h$. Dùng trong PCL (\verb|pcl::MovingLeastSquares|), Meshlab. Hạn chế: bandwidth $h$ — quá nhỏ giữ nhiễu, quá lớn làm mờ feature.

\textbf{Câu 23. WLOP là gì? Khác MLS thế nào?}\\
\textbf{WLOP} (Weighted Locally Optimal Projection, Huang 2009) sinh ra tập điểm \textbf{output} $\{x_i\}$ phân bố đều, đã khử nhiễu, từ input nhiễu $\{p_j\}$, ngay cả \textit{không cần} pháp tuyến. Tối thiểu hóa:
\[
E(X)=\sum_i\sum_j\theta(\|x_i-p_j\|)\|x_i-p_j\|-\mu\sum_i\sum_{i'\neq i}\theta(\|x_i-x_{i'}\|)\eta(\|x_i-x_{i'}\|)
\]
Số hạng đầu \textbf{hấp dẫn} kéo $x_i$ về dữ liệu input. Số hạng sau \textbf{đẩy} làm các điểm output trải đều. Trọng số $\theta$ giảm trọng vùng dày. \textbf{Khác MLS:} (a) MLS chiếu từng điểm input riêng và kế thừa không đều; WLOP phân bố lại; (b) WLOP không cần pháp tuyến; (c) MLS giả định bề mặt mượt, WLOP xử lý dữ liệu nhiều outlier tốt hơn. Dùng WLOP cho output scanner thô không đều cao.

\textbf{Câu 24. Ước lượng pháp tuyến point cloud bằng PCA?}\\
\textbf{Các bước:} 1) Với điểm $p_i$, tìm $k$-NN hoặc neighbors bán kính $r$ là $\{q_j\}$. 2) Tính centroid $\bar p$ và \textbf{ma trận hiệp phương sai} $C=\tfrac{1}{k}\sum(q_j-\bar p)(q_j-\bar p)^T$. 3) Phân rã trị riêng: $C=V\Lambda V^T$ với $\lambda_0\le\lambda_1\le\lambda_2$. Vector riêng ứng $\lambda_0$ nhỏ nhất là pháp tuyến $n_i$ (vuông góc mặt phẳng cục bộ). 4) Tỷ số $\lambda_0/(\lambda_0+\lambda_1+\lambda_2)$ là \textbf{độ biến thiên bề mặt}/proxy cho độ cong. 5) \textbf{Định hướng nhất quán}: dựng đồ thị Riemann của $k$-NN, xây cây bao trùm tối thiểu, lan truyền hướng bằng cách lật nếu $n_i\cdot n_j<0$ (Hoppe 1992). Biến thể robust: trọng số theo khoảng cách hoặc fit RANSAC.

\textbf{Câu 25. So sánh smoothing lưới và smoothing point cloud — khác nhau gì?}\\
\textbf{Smoothing lưới} dựa vào \textbf{kết nối} tường minh: mỗi đỉnh có one-ring cố định, nên Laplacian/Taubin/cotangent định nghĩa tốt và bộ lọc giữ feature có thể dùng góc nhị diện. \textbf{Smoothing point cloud} không có kết nối — phải xây neighborhood động bằng $k$-NN/bán kính mỗi lần. Phương pháp: MLS, WLOP, bilateral PC, hoặc deep learning. Smoothing lưới còn quản lý \textbf{topology} và đảm bảo manifold; PC không giữ topology. Lưới có pháp tuyến tường minh từ mặt; PC cần ước lượng PCA — bản thân nhiễu. Smoothing PC còn phải xử lý \textbf{biến thiên mật độ} (gần vs xa). Pipeline phổ biến: smooth PC $\to$ tái dựng lưới $\to$ smooth lưới.

\textbf{Câu 26. Mô tả pipeline geometric modeling đầy-đủ.}\\
\textbf{Các giai đoạn:} 1) \textbf{Thu nhận} — laser scan, photogrammetry, depth camera, hoặc CAD. 2) \textbf{Đăng ký} — gộp nhiều scan (ICP, SLAM toàn cục). 3) \textbf{Làm sạch} — xóa outlier (SOR/ROR), lấp khe. 4) \textbf{Downsample} — voxel grid, Poisson-disk. 5) \textbf{Ước lượng pháp tuyến} — PCA + định hướng nhất quán. 6) \textbf{Tái dựng bề mặt} — Poisson, Ball-Pivoting, Marching Cubes (từ implicit). 7) \textbf{Sửa lưới} — lấp lỗ, sửa cạnh non-manifold. 8) \textbf{Khử nhiễu} — Taubin, bilateral, MCF. 9) \textbf{Đơn giản hóa \& LOD} — QEM, progressive mesh. 10) \textbf{Parameterization \& UV} — cho texturing. 11) \textbf{Export} — OBJ/PLY/STL/glTF. 12) \textbf{Render / Ứng dụng} — game engine, in 3D, FEA, VR.

\textbf{Câu 27. Các định dạng 3D — so sánh OBJ, PLY, STL, OFF.}\\
\textbf{OBJ} (Wavefront): ASCII, hỗ trợ vertex, normal, UV, material (\verb|.mtl|), face đa giác (không chỉ tam giác). Chuẩn de-facto trao đổi asset. \textbf{PLY} (Stanford): ASCII hoặc binary, hỗ trợ thuộc tính tùy ý per-vertex/per-face (color, intensity, custom). Chuẩn cho output scanner và nghiên cứu. \textbf{STL}: ASCII hoặc binary, chỉ là tập tam giác với pháp tuyến mặt — \textit{không có index dùng chung, không có color, không có UV}. Chuẩn cho in 3D / CAD slicer. \textbf{OFF}: ASCII, rất đơn giản — header, vertex list, face list theo index. Dùng trong học thuật (Princeton Shape Benchmark). Cần biết thêm: \textbf{glTF 2.0} (binary, PBR, animation — web/games hiện đại), \textbf{FBX} (proprietary, animation), \textbf{USD} (Pixar, scene graph).

\textbf{Câu 28. Các biểu diễn 3D chính?}\\
1) \textbf{Point cloud} — tập điểm $(x,y,z)$ không thứ tự; native cho scanner; không topology. 2) \textbf{Lưới tam giác/đa giác} — đỉnh + cạnh + mặt; kết nối tường minh; chuẩn render. 3) \textbf{NURBS} — non-uniform rational B-spline; mượt, tham số, chính xác cho CAD. 4) \textbf{Implicit / SDF} — $f(x,y,z)=0$; CSG dễ, kiểm tra inside O(1); cần Marching Cubes để mesh. 5) \textbf{Voxel grid} — pixel 3D; native cho dữ liệu y khoa/khối; tốn bộ nhớ. 6) \textbf{CSG} — Constructive Solid Geometry, cây boolean của primitive; CAD cổ điển. 7) \textbf{Subdivision} — lưới thô + quy tắc tinh chỉnh; mượt tại mọi LOD. 8) \textbf{Neural fields} (NeRF, occupancy net) — MLP mã hóa hình học/radiance.

\textbf{Câu 29. Khi nào chọn lưới vs implicit vs voxel?}\\
\textbf{Lưới:} tốt nhất cho \textit{render} và \textit{game} — GPU yêu tam giác, bộ nhớ chỉ surface, UV tường minh. Kém với boolean và thay đổi topology. \textbf{Implicit / SDF:} tốt cho \textit{modeling}, \textit{boolean CSG}, \textit{ray-marching} (shadertoy), \textit{collision} — kiểm tra inside dạng đóng. Kém khi render trực tiếp không qua chuyển đổi. \textbf{Voxel:} tốt cho \textit{dữ liệu khối} (CT/MRI y khoa), \textit{mô phỏng} (chất lỏng, vật lý voxel), \textit{game kiểu Minecraft}. Kém về bộ nhớ: tăng $O(n^3)$ — giải bằng octree thưa (\textbf{VDB}, \textbf{NanoVDB}). \textbf{Nguyên tắc:} render = lưới, sculpt = SDF, simulate = voxel. Engine hiện đại trộn: Unreal Nanite dùng lưới, Lumen dùng SDF, mô phỏng chất lỏng dùng VDB.

\textbf{Câu 30. $C^0,C^1,C^2$ continuity là gì? Tại sao quan trọng?}\\
Tại biên giữa hai đường cong/bề mặt: \textbf{$C^0$} = vị trí trùng (không có khe). \textbf{$C^1$} = vị trí và đạo hàm bậc 1 (\textbf{tiếp tuyến}) trùng — bề mặt mượt thị giác, không có ``gãy''. \textbf{$C^2$} = vị trí, tiếp tuyến, đạo hàm bậc 2 (\textbf{độ cong}) trùng — mượt hơn nữa, quan trọng cho highlight và phản xạ. \textbf{$G^n$} (continuity hình học) là phiên bản lỏng hơn: \textit{hướng} tiếp tuyến/độ cong trùng, độ lớn có thể khác. \textbf{Tại sao quan trọng:} 1) Class-A surface ô tô/sản phẩm cần $G^2$ cho phản xạ sạch. 2) Đường animation cần ít nhất $C^1$ để tránh giật. 3) Subdivision (Catmull-Clark) cho $C^2$ trừ đỉnh đặc biệt. 4) Bezier patch cần khớp tiếp tuyến biên cẩn thận để có $C^1$.

\textbf{Câu 31. Khử nhiễu point cloud trong LiDAR / xe tự lái?}\\
LiDAR có nhiễu và điểm giả từ \textbf{mưa, sương, tuyết, khói thải, bụi, multipath sensor}. Pipeline trong AV (Waymo, Tesla, Mobileye): 1) \textbf{Lọc range} — xóa return rất gần/rất xa. 2) \textbf{ROR/SOR} — xóa điểm trôi đơn (giọt mưa). 3) \textbf{Khớp mặt đất} (RANSAC) — tách mặt đất khỏi vật. 4) \textbf{Voxel downsample} — cho 3D detector real-time (PointPillars, CenterPoint). 5) \textbf{Khử nhiễu thời gian} — tích lũy frame, vote bền vững (điểm mưa ngẫu nhiên không lặp lại). 6) \textbf{Lọc theo thời tiết} — DROR (Dynamic ROR) cho tuyết. Cloud sạch là then chốt cho \textbf{phát hiện vật 3D}, \textbf{tracking}, \textbf{SLAM}. False positive nhỏ cũng có thể gây phanh ảo.

\textbf{Câu 32. Đơn giản hóa lưới trong render real-time / game?}\\
Game không thể render asset 100M tam giác ở 60-120 Hz. \textbf{Ứng dụng:} 1) \textbf{Chuỗi LOD rời rạc} (LOD0-4) sinh bởi QEM offline, chuyển theo kích thước màn hình. 2) \textbf{HLOD} — cụm tòa nhà xa gộp thành proxy đơn (Unreal). 3) \textbf{Imposter} — đơn giản hơn cả lưới, chỉ là billboard từ LOD cực thấp. 4) \textbf{Shadow LOD} — đơn giản hóa aggressive hơn riêng cho shadow caster. 5) \textbf{Collision LOD} — lưới vật lý đơn giản runtime. 6) \textbf{Streaming} — progressive mesh / Nanite virtualized geometry. 7) \textbf{VR/mobile} — LOD aggressive vì ngân sách nhiệt/perf chặt. Tổ hợp tiết kiệm \textbf{vertex shader cost}, draw call, băng thông, shadow pass. Không có simplification, game open-world hiện đại không khả thi.

\textbf{Câu 33. Khử nhiễu trong scan 3D y khoa?}\\
CT/MRI/siêu âm và scanner bề mặt (intra-oral nha khoa, body scanner) cho khối/lưới nhiễu. \textbf{Ứng dụng:} 1) \textbf{Tái dựng CT} — thuật toán lặp (SART, MBIR) gồm prior khử nhiễu (TV, BM3D). 2) \textbf{Trích bề mặt} — Marching Cubes rồi Taubin smoothing cho bề mặt nội tạng. 3) \textbf{Thiết kế mão răng} — bilateral mesh denoising giữ feature đỉnh/rãnh trong khi xóa speckle scanner. 4) \textbf{Lập kế hoạch phẫu thuật} — khử nhiễu point cloud trên scan intra-op đã đăng ký để overlay AR. 5) \textbf{Khớp implant/prosthesis} — lưới sạch để in 3D. 6) \textbf{Radiomics / AI} — ảnh đã khử nhiễu cho feature đáng tin hơn. Khử nhiễu phải \textbf{giữ feature} — không được làm mờ biên khối u hoặc thành mạch, vì vậy bilateral và edge-aware chiếm ưu thế.

\textbf{Câu 34. Tại sao geometric modeling là nền tảng của digital twin?}\\
\textbf{Digital twin} là bản sao 3D ảo đồng bộ của tài sản vật lý (nhà máy, tòa nhà, thành phố, cơ thể). Geometric modeling cung cấp: 1) \textbf{Thu nhận} — laser scan / photogrammetry để chụp thực tế as-built. 2) \textbf{Làm sạch + khử nhiễu} — mô hình dùng được từ scan thô. 3) \textbf{Đơn giản hóa + LOD} — twin thành phố có hàng tỷ tam giác; LOD mới dùng tương tác được. 4) \textbf{Lưới + parameterization} — đặt sensor, mô phỏng (CFD, FEM), phân tích nhiệt. 5) \textbf{Cập nhật liên tục} — re-scan, đăng ký, diff với lưới trước để phát hiện thay đổi (vd biến dạng kết cấu). 6) \textbf{Tương tác liên thông} — định dạng chuẩn (USD, IFC, glTF) chia sẻ giữa công cụ. Không có geometric modeling, digital twin chỉ là CAD drawing — thiếu tính trung thực cao, sensor-driven, tiến hóa thời gian làm nên ``twin''.

\textbf{Câu 35. QEM kết hợp Marching Cubes trong y khoa thế nào?}\\
\textbf{Pipeline:} 1) Thu khối CT/MRI. 2) Phân đoạn nội tạng bằng ngưỡng/AI để có trường nhị phân hoặc vô hướng. 3) Chạy \textbf{Marching Cubes} (Lorensen 1987) trích iso-surface — thường \textit{vài triệu tam giác nhỏ}, có aliasing răng cưa. 4) Áp dụng \textbf{Taubin smoothing} để xóa artifact răng cưa voxel. 5) Chạy \textbf{QEM simplification} để giảm còn vài chục nghìn tam giác, giữ feature giải phẫu (cong mạch, hình khối u). 6) Tính \textbf{pháp tuyến + UV} cho shading. 7) Render PBR cho lập kế hoạch phẫu thuật, overlay VR/AR, hoặc in 3D. Tổ hợp này thiết yếu vì output Marching Cubes quá dày cho VR tương tác nhưng có topology đúng — QEM + smoothing biến chúng thành thực dụng, còn QEM giữ feature giữ chi tiết quan trọng lâm sàng (vd nhánh mạch, biên khối u).

\end{paracol}
